\section{Related Work and Scope}
\paragraph{Cellwise detection and imputation.}
DDC uses robust relationships between numerical variables to predict cells and assess
their residuals~\cite{rousseeuw2018ddc}. Its published algorithm includes robust
standardization, bivariate relations, prediction and deshrinkage. Our numerical
DDCcore baseline is a simpler comparator and cannot establish a ranking against
reference DDC. CellGMM addresses numerical populations with a mixture model
and a detection--imputation step~\cite{zaccaria2025cellgmm}; discarding an entire
record can lose useful information from its other cells. Our removal operates
on whole context rows for one target at a time and offers no corresponding
mixture-estimation or cell-repair guarantee. Cellwise trimming has also been built
into model-based co-clustering~\cite{fibbi2026celltricc}, so trimming as such is not new here. Conformal Data Cleaning (CDC) combines column prediction with
calibrated prediction sets or intervals~\cite{jager2024cdc}. Its evaluated deployment
scenario trains on clean data and subsequently encounters corrupted inputs.
Our deployment arms receive one dirty table and no clean calibration split.
The LightGBM conditional comparator is inspired by column prediction but does not
reproduce CDC, and our scores have no conformal coverage interpretation.

\paragraph{Foundation models and contaminated contexts.}
TabPFN provides pretrained conditional prediction for small tabular
problems~\cite{hollmann2025tabpfn}. We use pinned second-generation weights through
an explicit adapter, rather than whichever model a current package would select
by default. Unsupervised extensions of TabPFN also use conditional prediction;
we do not claim this reduction as novel or call our raw arm a reproduction of
the extensions. TACTIC explicitly studies contaminated contexts for row anomaly
detection~\cite{tactic2026}. Its anomaly-oriented pretraining and row-level
evaluation address a related robustness problem. Consequently, context contamination
itself is established prior work; our intervention concerns cell-level rankings
under fixed pretrained weights and nested selection, not a row-AUROC comparison.

\paragraph{Data quality and context selection for TFMs.}
QualiTab evaluates imperfect data and contrasts perfect-data, perfect-context
and no-intervention settings for TabPFN embeddings~\cite{papastergios2025wild}.
Its clean-context comparison is close to the motivation of our oracle diagnostic;
the benefit of cleaner contexts is therefore not a new hypothesis introduced here.
Model-aware cleaning learns reinforcement-learning policies that sequence cleaning
operators for TFM-oriented data
preparation~\cite{bertiequille2026prior}. We instead evaluate one fixed removal
rule as a cell-ranking intervention without learning a cleaning policy.
LUCoS selects representative contexts using unsupervised PFN embeddings,
then obtains labels for supervised prediction~\cite{ipas2026lucos}.
Our contexts contain observed target values from the dirty table, with no
annotation oracle for deployment. A study of test-time compute for TFMs reports that
attention-guided row retrieval helps on some large tables but can fail badly, for
example by distorting class composition, and that context expansion gives no
consistent gain~\cite{ning2026ttc}. It concerns supervised prediction from labelled
contexts, not detecting errors in a dirty context, so it is related evidence rather
than a direct comparator. Together these works rule out broad novelty claims about
TFM data cleaning, label-free context selection or context selection in general.

\paragraph{Systems for heterogeneous error detection.}
Raha is a configuration-free error-detection system and supplies the benchmark
repository used here~\cite{mahdavi2019raha}.
Auto-Test learns semantic-domain constraints from a large table corpus offline and
applies them to detect erroneous column values without per-table
supervision~\cite{chen2025autotest}. MaTElDa is semi-supervised: it groups cells of
several tables by domain and quality similarity so that a small labelling budget
can be shared across tables~\cite{ahmadi2025matelda}. LUCARIO targets unsupervised,
lightweight detection of pattern anomalies and outliers~\cite{qin2026lucario}.
These are direct cell-level competitors with different knowledge sources (a table
corpus, user labels, inferred constraints); none was run under our protocol. We use its paired dirty/clean tables,
not a reimplementation of its complete detection workflow. ZeroED combines
LLM-derived representations and labels with detector training~\cite{zeroed2025}.
That information and compute regime differs from our label-free local deployment
arms.
Neither system's published headline score is an interchangeable baseline result
for our capped-table protocol.

\paragraph{Measured comparators.}
We run frequency and pattern scores, a robust numerical score, column-wise
Isolation Forest~\cite{liu2008isolation}, DDCcore and conditional
LightGBM~\cite{ke2017lightgbm}. They test whether marginal rarity, simple numerical
deviation or the choice of conditional model explains an effect; they cannot rank
the declared rule against every context-selection strategy.
